\section{Further theory}\label{app:theory}

\subsection{The null distribution at a fixed threshold}\label{app:nullmean}

Call a feature null if it is continuous and independent of $Y$. Given the labels, the ranks of the $\nmin$ positives among the $n$ values of a null feature are a uniformly random subset, so a split that sends a fixed number $n_L$ of samples left has $k_L\sim\mathrm{Hypergeometric}(n,\nmin,n_L)$.

\begin{proposition}[exact null mean]\label{prop:nullmean}
At the root, for a null feature, a fixed (data-independent) $n_L$, the prevalence intercept and $\lambda = 0$, $\E[\Gain] = \E[X^2] = n/(n-1)$ for every $n_L$ and every $\nmin$.
\end{proposition}

\begin{proof}
With $b = \phat$, $G = 0$, so $G_R = -G_L$ with $G_L = n_L\phat - k_L$, and $\Gain = G_L^2(1/H_L + 1/H_R) = G_L^2\,n/(\phat(1-\phat)n_Ln_R)$. Here $\E G_L = 0$ and $\operatorname{Var}k_L = n_L\phat(1-\phat)(n-n_L)/(n-1)$; substituting gives $n/(n-1)$.
\end{proof}

Imbalance therefore does not change the expected gain of a single uninformative split; it changes the \emph{shape} of its distribution. $X^2$ is approximately $\chi^2_1$ only when both children are expected to hold many positives. For a child of size $m\ll n$ holding $j$ positives, $X^2\approx(j - m\phat)^2/(m\phat)$: with $m\phat = 0.1$ and $j = 2$ it is 36, against a 99.9th percentile of 10.8 for $\chi^2_1$.

\begin{proposition}[\code{min\_child\_weight} counts expected positives]\label{prop:mcw}
In the first tree, $H_L\ge w$ means $n_L\ge w/(b(1-b))$. Under the null, a child of that size holds on average $n_L\phat\ge w\phat/(b(1-b))$ positives. With the prevalence intercept this is $w/(1-\phat)\approx w$; with the XGBoost 2.x default at $\phat = 0.001$ ($b\approx 0.12$) it is about $0.0095\,w$, a child of about 10 samples.
\end{proposition}

So XGBoost's default $w = 1$ rules out, under the prevalence intercept, the splits whose children are expected to hold fewer than about one positive, which is where the heavy tail lives; under the 2.x default it protects almost nothing. With \code{hist}, children are also unions of bins, which in practice bounds them below by roughly $n/\code{max\_bin}$ samples; this is a description of the binning, not part of the proposition.

\subsection{Signal at a fixed threshold}\label{app:fixed}

\begin{proposition}[signal at a fixed threshold]\label{prop:fixed}
In the limit of \cref{thm:limit}, $N(u) \sim \Bin(\nmin, G(u))$, and the standardised deviation $(N(u)-\nmin u)^2/(\nmin u(1-u))$, which is $\Gamma(u)$ for $\lambda = 0$, has mean $G(u)(1-G(u))/(u(1-u)) + \nu(u)$ with non-centrality $\nu(u) = \nmin (G(u)-u)^2/(u(1-u))$. If $\nmin \to \infty$ along a sequence of placement laws $G_{\nmin}$ with $\sqrt{\nmin}\,(G_{\nmin}(u) - u) \to -c$ for a fixed $u$, the standardised count $(N(u) - \nmin u)/\sqrt{\nmin u(1-u)}$ converges in law to $\mathcal N(-c/\sqrt{u(1-u)},\,1)$. In the Gaussian model,
\begin{equation}\label{eq:twopi}
  \lim_{\delta\to 0}\, \frac{\max_u \nu(u)}{\nmin\delta^2} = \frac{2}{\pi},
\end{equation}
attained at the median split $u \to 1/2$.
\end{proposition}

The proof is in \cref{app:proof-fixed}. The ratio $\varphi(x)^2/(\Phi(x)(1-\Phi(x)))$ that appears is the efficiency of dichotomising a normal sample at $x$ \citep{cox1957note}. At finite $\delta$ the best threshold moves into the positives' tail and the maximum exceeds $(2/\pi)\nmin\delta^2$, by a factor of 1.03, 1.09 and 1.47 at $\delta = 0.3$, 0.5 and 1. So $(2/\pi)s$ is the small-effect value, not an upper bound.

\subsection{Efficiency of one split relative to the rank test}\label{app:efficiency}

\begin{corollary}[efficiency of one split relative to the rank test]\label{cor:efficiency}
Let a feature follow a location family with bounded standardised density $f$, shift $\delta \to 0$ and $\nmin\delta^2 \to s$. A pre-specified threshold $u$, with $F^{-1}(u)$ uniquely defined and $f$ continuous there, has non-centrality $s\,J_f(u)^2/(u(1-u))$ (\cref{prop:fixed}), and the rank test has non-centrality $12 s\,(\int f^2)^2$ (\cref{prop:mw}, since $\AUC - \tfrac12 = \delta\int f^2 + o(\delta)$). The best fixed threshold, chosen knowing $f$, therefore has asymptotic relative efficiency
\begin{equation}\label{eq:efficiency}
  \mathrm{ARE}(f) = \frac{\sup_u J_f(u)^2/(u(1-u))}{12\,\bigl(\int f^2\bigr)^2}
\end{equation}
with respect to the rank test: $2/3$ for the normal, $3/4$ for the logistic, $0.85$ for $t_3$ and $4/3$ for the Laplace distribution (all at the median), and $4.8$ for the log-normal with $\sigma = 1$, at an interior maximiser $u \approx 0.01$ (the median split has $0.40$).
\end{corollary}

For the normal, $2/3$ is the Pitman efficiency of the median test relative to the Wilcoxon test \citep{mood1954asymptotic,hodges1956efficiency}: one fixed split needs 1.5 times as many positives as the rank test. For heavy-tailed features the gap narrows, and for strongly skewed ones a well-placed split beats the rank test; the log-normal value needs $\nmin u^* \gg 1$, so it matters only for $\nmin$ well above 100.

\subsection{Histogram search}\label{app:hist}

\begin{corollary}[histogram search]\label{cor:hist}
With \code{tree\_method="hist"} and a fixed \code{max\_bin} $= B$, suppose (H1) the cut points have empirical-CDF positions converging to $j/B$ almost surely, and (H2) no $j/B$ equals $w/\nmin$ or $1-w/\nmin$. Then the conclusion of \cref{thm:limit} holds with $u$ restricted to $\{1/B,\dots,(B-1)/B\}\cap[w/\nmin, 1-w/\nmin]$.
\end{corollary}

\begin{proof}
Under (H1) the cut $c^n_j$ has $F_n(c^n_j)\to j/B$, where $F_n$ is the empirical distribution function of all $n$ values, so $k_L$ at that cut is $N_n(F_n(c^n_j)) = N(\ell_n(F_n(c^n_j)))\to N(j/B)$, because $N$ is almost surely continuous at $j/B$. The bracket converges as before. Under (H2) each cut is eventually admissible if and only if $j/B\in(w/\nmin, 1-w/\nmin)$, so the maximum over cuts converges to the maximum of $\Gamma$ over the grid.
\end{proof}

(H1) is an assumption about XGBoost's quantile sketch. In the first tree the Hessian is constant, so the sketch is unweighted, but we have not verified (H1) analytically. Every prediction for \code{hist} in this paper is therefore conditional on it.

\subsection{The full threshold window}\label{app:conj1}

\begin{conjecture}[slow drift on the full window]\label{conj:drift}
Let $M^{(0)}_{\nmin}$ be the null supremum of \eqref{eq:limit} on the full window $[w/\nmin, 1-w/\nmin]$, with $\lambda, w > 0$ fixed, and let $\omega(s,K;\nmin)$ be the root win probability of an informative Gaussian feature with $\nmin\delta^2 = s$ against $K$ null features on that window. Then (a) $M^{(0)}_{\nmin}/(2\log\log\nmin) \to 1$ in probability as $\nmin\to\infty$; (b) at fixed $s$ and $K$, $\omega(s,K;\nmin) \to 1/(K+1)$; (c) in finite range, $\omega(12.5, 9;\nmin)$ is non-increasing in $\nmin$ and falls below 0.75 by $\nmin = 10^5$.
\end{conjecture}

The unregularised analogue: the squared supremum of the standardised uniform empirical process over $[a_m, 1-a_m]$, with $a_m \to 0$ and $m a_m \to \infty$, is $2\log\log(1/a_m)(1 + o_P(1))$ \citep{jaeschke1979asymptotic,csorgo1993weighted}. What is unproved is the regularised boundary region, where $\nmin a_{\nmin} = w$ stays bounded and the classical result does not apply. Matching the centre drift $2s/\pi$ against the null level $2\log\log\nmin$ suggests, as a heuristic, that the win probability tends to $1/(K+1)$ for $s = c\log\log\nmin$ with $c < \pi$ and to 1 with $c > \pi$, so detection at fixed power needs $s$ to grow like $\log\log\nmin$. Over the range we can compute, the median null supremum is about $1.34\times 2\log\log\nmin$, so part (a) is an asymptotic statement that the evidence does not yet reach.


Numerical evaluation of the limit \eqref{eq:limit} with $\lambda = w = 1$ and $K = 9$ null features (1{,}500 signal and 6{,}000 null draws per column; Monte Carlo standard error about 0.01):

\begin{center}\small
\begin{tabular}{lcccccc}
\toprule
$\nmin$ & 10 & 50 & 200 & 1{,}000 & 5{,}000 & 20{,}000\\
\midrule
median null supremum & 2.40 & 3.63 & 4.54 & 5.24 & 5.74 & 6.14\\
$\P$(informative wins), $s = 12.5$, full window & 0.87 & 0.85 & 0.82 & 0.79 & 0.80 & 0.76\\
\bottomrule
\end{tabular}
\end{center}

The median grows roughly like $2\log\log\nmin$ and the win probability decreases slowly; on the fixed window $[0.1, 0.9]$ the win probability is flat at $0.86\pm0.01$ for $\nmin$ from 50 to 5{,}000. The library computes these probabilities by Monte Carlo on \eqref{eq:limit} (\code{theory.detection\_probability}).

\subsection{Gain accounting for the full ensemble}\label{app:accounting}

\begin{proposition}[gain accounting]\label{prop:accounting}
Consider one tree with $\gamma = 0$ and leaves $\ell$ with gradient and Hessian sums $(G_\ell, H_\ell)$, and let $(G,H)$ be the sums over all samples. The reported gains of the tree's splits add up to
\begin{equation}\label{eq:telescope}
  \sum_{\text{splits}}\Gain = \sum_\ell\frac{G_\ell^2}{H_\ell+\lambda} - \frac{G^2}{H+\lambda}.
\end{equation}
For this tree, given its structure, with leaf values $-\eta G_\ell/(H_\ell+\lambda)$, the decrease of the second-order (Newton) approximation of the training loss is
\begin{equation}\label{eq:newton}
  \Delta\tilde L = \eta\Bigl(1-\frac\eta2\Bigr)\sum_\ell\frac{G_\ell^2}{H_\ell+\lambda} + \frac{\eta^2\lambda}{2}\sum_\ell\frac{G_\ell^2}{(H_\ell+\lambda)^2}.
\end{equation}
\end{proposition}

\begin{proof}
Each split's gain is children minus parent. Summing over internal nodes, every internal node other than the root appears once as a child and once as a parent, so the sum telescopes to \eqref{eq:telescope}. For \eqref{eq:newton}, the Newton model of the loss change is $\sum_\ell(G_\ell v_\ell + \tfrac12H_\ell v_\ell^2)$ with $v_\ell = -\eta G_\ell/(H_\ell+\lambda)$; write $H_\ell = (H_\ell+\lambda)-\lambda$ and collect terms.
\end{proof}

\emph{Heuristic consequence.} When $\lambda\ll H_\ell$ and $G\approx 0$, a tree's total gain is its Newton loss decrease divided by $\eta(1-\eta/2)$, which is 10.5 at $\eta = 0.1$ (the neglected $\lambda$ term is at most $\eta/(2-\eta)\approx 5\%$). If the Newton model approximates the true loss change, the \code{total\_gain} recorded by all features over the ensemble is about $10.5\times$ the training log-loss removed. That loss starts at $nH_2(\phat)\approx\nmin(1+\ln(1/\phat))$ nats under the prevalence intercept, where $H_2$ is the binary entropy, so more loss is available to memorise at small $\phat$. The identity concerns \code{total\_gain}; the selection statistic in most of this paper is the mean \code{gain}, and we have no argument linking the two. The measured fraction of the initial loss removed is a Newton-approximation estimate (\cref{app:agreement}, row P15).

\subsection{Keep rate of a null candidate in Boruta-style selection}\label{app:rounds-grid}

\Cref{tab:rounds-grid} evaluates $\P(\Bin(R,\pi) > r_R)$ for a constant hit rate $\pi = 1/(K+1)$, the keep rate \cref{prop:rounds} would give if the hit rate did not vary between data sets. It grows with the number of candidates $d$, because the Bonferroni threshold $r_R$ falls, and it is not monotone in $R$.

\begin{table}[h]
\centering\small
\caption{Keep rate of a null candidate under a constant hit rate $1/(K+1)$, with $K = 10$ shadows, $\alpha = 0.05$ and a Bonferroni correction over $d$ candidates, as a function of the number of rounds $R$. In parentheses: the rejection threshold $r_R$.}\label{tab:rounds-grid}
\begin{tabular}{lcccccc}
\toprule
$d$ & $R = 10$ & 18 & 20 & 40 & 60 & 100\\
\midrule
10 & 0.614 (0) & 0.074 (3) & 0.102 (3) & 0.000 (11) & 0.000 (19) & 0.000 (36)\\
50 & 0.614 (0) & 0.221 (2) & 0.272 (2) & 0.003 (9) & 0.000 (17) & 0.000 (34)\\
200 & 1.000 ($-1$) & 0.496 (1) & 0.272 (2) & 0.009 (8) & 0.000 (16) & 0.000 (32)\\
1000 & 1.000 ($-1$) & 0.820 (0) & 0.554 (1) & 0.026 (7) & 0.000 (14) & 0.000 (30)\\
\bottomrule
\end{tabular}
\end{table}

In practice the hit rate varies between data sets, and the measured keep rate at large $R$ is set by that variation (\cref{sec:methods}); the table gives the part of the keep rate that is due to the test alone.
